\worksheettitle
  {M02\(\star\). Где находится второй конец?}
  {\modulemark{M02\(\star\)} Усложнение}

\studentnameblock

\begin{notebox}[Идея]
  При измерении по двум отметкам находим длину вычитанием. В обратной задаче
  известны одна отметка и длина: второй конец можно получить сложением или
  вычитанием.
\end{notebox}

\begin{practicebox}[1. Движение вправо]
  Точка $A$ находится у отметки $2\text{ см }4\text{ мм}$. Отрезок
  $AB$ длиной $3\text{ см }7\text{ мм}$ отложили вправо.
  У какой отметки находится $B$? \writinglines{1}
\end{practicebox}

\begin{practicebox}[2. Два возможных направления]
  Точка $C$ находится у отметки $6\text{ см }5\text{ мм}$. Известно, что
  $CD=2\text{ см }8\text{ мм}$. Найдите обе возможные отметки точки $D$,
  если отрезок можно отложить влево или вправо.
  \writinglines{2}
\end{practicebox}

\begin{practicebox}[3. Ограничение линейки]
  На шкале от $0$ до $10\text{ см}$ точка $M$ стоит у $1\text{ см}$.
  Нужно отложить $MN=2\text{ см }6\text{ мм}$. Какие из двух направлений
  допустимы на этой шкале? Объясните. \writinglines{2}
\end{practicebox}

\begin{checkpointbox}[Создайте задачу]
  Придумайте положение первого конца и длину так, чтобы второй конец мог иметь
  две допустимые отметки на шкале от $0$ до $10\text{ см}$.
  \writinglines{2}
\end{checkpointbox}
